\begin{table}[htbp]\centering\small
\caption{같은 양의 수준 예측에 대한 손실함수별 개선율}\label{tab:common_loss}
\begin{tabular}{lrrrrr}\toprule
종목 & $\beta$ & $h$ & 수준 MSE & 로그 MSE & QLIKE\\\midrule
\multicolumn{6}{l}{\textit{상장 전}}\\
삼성전자 & $+2$ & 5 & 3.33 & 2.48 & 3.18\\
삼성전자 & $-2$ & 5 & 3.02 & 2.42 & 3.00\\
삼성전자 & $+2$ & 20 & 7.33 & 5.28 & 5.77\\
삼성전자 & $-2$ & 20 & 7.18 & 5.10 & 5.49\\
SK하이닉스 & $+2$ & 5 & 2.41 & -0.37 & 1.49\\
SK하이닉스 & $-2$ & 5 & 2.30 & -0.40 & 1.44\\
SK하이닉스 & $+2$ & 20 & 2.78 & 0.71 & 1.78\\
SK하이닉스 & $-2$ & 20 & 2.33 & 0.40 & 1.45\\
\midrule
\multicolumn{6}{l}{\textit{상장 후}}\\
삼성전자 & $+2$ & 5 & -2.18 & -2.94 & -14.03\\
삼성전자 & $-2$ & 5 & -3.00 & -3.98 & -24.84\\
삼성전자 & $+2$ & 20 & 2.33 & -3.65 & -15.01\\
삼성전자 & $-2$ & 20 & -1.06 & -6.23 & -22.56\\
SK하이닉스 & $+2$ & 5 & -0.72 & 2.35 & -4.85\\
SK하이닉스 & $-2$ & 5 & -1.80 & 1.71 & -12.07\\
SK하이닉스 & $+2$ & 20 & 3.93 & 2.90 & -1.65\\
SK하이닉스 & $-2$ & 20 & 1.38 & 0.83 & -7.67\\
\bottomrule\end{tabular}
\smallskip
\begin{minipage}{0.97\textwidth}\footnotesize 단위는 \%. 기존 로그 설명변수 사양의 M1과 M2를 각각 훈련잔차 기반으로 재변환한 후, 같은 수준 $G$와 같은 양의 예측값으로 세 손실을 계산하였다. 단순 로그회귀 예측 자체를 평가한 기존 표와 대상 예측값이 다르다. 두 모형 모두 이 표의 표본에서 양수 하한에 걸린 예측은 없다.\end{minipage}
\end{table}
